\documentclass[aspectratio=169,11pt]{beamer}

\usepackage[T1]{fontenc}
\usepackage[utf8]{inputenc}
\usepackage[english]{babel}
\usepackage{graphicx}
\usepackage{booktabs}
\usepackage{amsmath}
\usepackage{hyperref}

\usetheme{default}
\setbeamertemplate{navigation symbols}{}
\setbeamertemplate{footline}[frame number]

\definecolor{ustnavy}{HTML}{0B1F3A}
\definecolor{ustteal}{HTML}{1F6F8B}
\definecolor{ustaccent}{HTML}{8A2B0E}
\definecolor{ustgray}{HTML}{4A5560}
\definecolor{ustgreen}{HTML}{1B5E3B}

\setbeamercolor{structure}{fg=ustteal}
\setbeamercolor{frametitle}{fg=ustnavy}
\setbeamercolor{title}{fg=ustnavy}
\setbeamercolor{subtitle}{fg=ustgray}
\setbeamercolor{author}{fg=ustgray}
\setbeamercolor{institute}{fg=ustgray}
\setbeamercolor{item}{fg=ustteal}
\setbeamercolor{block title}{fg=white,bg=ustteal}
\setbeamercolor{block body}{fg=ustnavy,bg=ustnavy!4}
\setbeamercolor{block title alerted}{fg=white,bg=ustaccent}
\setbeamercolor{block body alerted}{fg=ustnavy,bg=ustaccent!8}
\setbeamercolor{block title example}{fg=white,bg=ustgreen}
\setbeamercolor{block body example}{fg=ustnavy,bg=ustgreen!8}
\setbeamerfont{title}{series=\bfseries,size=\Large}
\setbeamerfont{frametitle}{series=\bfseries,size=\large}

\newcommand{\shot}[2]{%
  \IfFileExists{#1}{%
    \includegraphics[width=#2\textwidth]{#1}%
  }{%
    \fbox{\begin{minipage}{0.88\textwidth}\centering\vspace{1.1cm}
    \textbf{Image expected}\\[0.35em]
    {\small\ttfamily\detokenize{#1}}\\[0.35em]
    {\small Please add the file to \texttt{user\_shots/}.}
    \vspace{1.1cm}\end{minipage}}%
  }%
}

\title{Universal Spectrum Toolkit}
\subtitle{Conversion of MCA spectrum files into ROOT analysis objects}
\author{Eren Yılmaz}
\institute{Technical presentation}
\date{2026}

\begin{document}

%---------------------------------------------------------------------
\begin{frame}[plain]
  \titlepage
  \vspace{-0.4em}
  \begin{center}
    \small\textcolor{ustgray}{Problem definition, methodological approach, implementation, and usage}
  \end{center}
\end{frame}

%---------------------------------------------------------------------
\begin{frame}{Outline}
  \begin{enumerate}
    \item Problem definition: why ROOT cannot process \texttt{.spe} files as spectra
    \item Incorrect reading: the result of interpreting data without a known layout
    \item Method: the logic constructed to resolve the binary structure
    \item Mathematical decision: scoring of layout hypotheses
    \item Application output and validation
    \item Usage: menu-based interface
  \end{enumerate}
\end{frame}

%---------------------------------------------------------------------
\begin{frame}{Problem definition}
  \begin{alertblock}{Core problem}
    ROOT is an analysis framework. Its expected data model consists of
    \texttt{.root} files, \texttt{TH1}, \texttt{TTree}, and built-in I/O definitions.
    MCA-produced \texttt{.spe}/\texttt{.chn} files, by contrast, are raw byte sequences.
  \end{alertblock}

  \vspace{0.5em}
  \textbf{Consequence:}
  \begin{itemize}
    \item Physically opening a file does not imply correctly reading a spectrum.
    \item \texttt{.spe} is not a single international standard.
    \item Header size, numeric type, and byte order vary by vendor.
  \end{itemize}
\end{frame}

%---------------------------------------------------------------------
\begin{frame}{File content without the macro}
  \begin{center}
    \shot{user_shots/03_blind_spe.png}{0.78}
  \end{center}
\end{frame}

%---------------------------------------------------------------------
\begin{frame}{What this evidence demonstrates}
  \begin{itemize}
    \item Opening the file as text does not produce a readable spectrum.
    \item The hex dump reveals only a limited header signature
          (\texttt{DataR\_CH}).
    \item Ortec ASCII tags (\texttt{\$DATA}, \texttt{\$ENER\_FIT}, \ldots)
          are absent; therefore text-based reading is invalid.
  \end{itemize}

  \vspace{0.7em}
  \begin{block}{Conclusion}
    The problem is not printing the file to the screen. The problem is
    assigning the correct \textbf{semantic mapping} to the byte sequence:
    header offset, numeric type, channel count, and the counts array.
  \end{block}
\end{frame}

%---------------------------------------------------------------------
\begin{frame}{Methodological idea}
  \begin{exampleblock}{Central approach}
    Rather than interpreting the file directly as a spectrum, first resolve
    the binary layout through systematic hypotheses; then encode the
    validated layout in software as an explicit rule.
  \end{exampleblock}

  \vspace{0.6em}
  \textbf{Three-step reasoning:}
  \begin{enumerate}
    \item Inspect the raw file at the byte level
    \item Generate and score candidate layout hypotheses
    \item Convert the winning layout into a channel--counts array, then into a ROOT object
  \end{enumerate}
\end{frame}

%---------------------------------------------------------------------
\begin{frame}{Resolving the binary structure}
  \begin{center}
    \shot{user_shots/04_binary_probe.png}{0.90}
  \end{center}
\end{frame}

%---------------------------------------------------------------------
\begin{frame}{Mathematical decision model}
  Given a file size $S$, each hypothesis is defined by:

  \[
  H = S - N\cdot B
  \]

  \begin{itemize}
    \item $N$: assumed channel count (e.g.\ $16384$)
    \item $B$: bytes per value (e.g.\ $B=4$ for \texttt{float32})
    \item $H$: header offset (e.g.\ $H=40 \Rightarrow S=65576$)
  \end{itemize}

  \vspace{0.3em}
  The payload is decoded with the selected type and endianness; candidates are scored by:
  \begin{itemize}
    \item \texttt{finite}: fraction of finite values
    \item \texttt{negative}: fraction of negative values (unexpected in spectra)
    \item \texttt{smoothness}: physical consistency of the sequence
    \item \texttt{score}: composite decision metric
  \end{itemize}
\end{frame}

%---------------------------------------------------------------------
\begin{frame}{Transferring this logic into code}
  \begin{block}{Discovery stage}
    The probe output shows that the highest-scoring layout for the sample data is\\
    $H=40$, $N=16384$, \texttt{float32}.
  \end{block}

  \vspace{0.5em}
  \begin{exampleblock}{Production rule}
    This layout is not left as a random guess.\\
    It is converted into a validated reader rule in code:\\
    header + channel count + numeric type + byte order.
  \end{exampleblock}

  \vspace{0.5em}
  \textbf{Result:} The software no longer reads the file by speculation,\\
  but through a previously validated semantic model.
\end{frame}

%---------------------------------------------------------------------
\begin{frame}{Validation of the conversion}
  \begin{center}
    \shot{user_shots/02_toolkit_match.png}{0.90}
  \end{center}
\end{frame}

%---------------------------------------------------------------------
\begin{frame}{Meaning of the decision output}
  \begin{itemize}
    \item The ASCII reader does not match: the file is not text.
    \item The GF3-like reader matches: $40$-byte header + \texttt{float32} LE.
    \item Channel count: $16384$
    \item Produced objects:
      \begin{itemize}
        \item \texttt{TH1D hSpectrumChannel}
        \item \texttt{TTree spectrum}
        \item \texttt{conversion\_metadata}
        \item raw-byte archive
      \end{itemize}
  \end{itemize}

  \vspace{0.5em}
  \begin{block}{Summary}
    The raw byte sequence has been converted, in a controlled manner,
    into a spectrum object that ROOT can analyze.
  \end{block}
\end{frame}

%---------------------------------------------------------------------
\begin{frame}{Successful conversion output}
  \begin{columns}[c]
    \begin{column}{0.48\textwidth}
      \centering
      \shot{user_shots/05_converted_linear.png}{1.0}\\
      \scriptsize Linear scale
    \end{column}
    \begin{column}{0.48\textwidth}
      \centering
      \shot{user_shots/06_converted_log.png}{1.0}\\
      \scriptsize Logarithmic scale
    \end{column}
  \end{columns}
\end{frame}

%---------------------------------------------------------------------
\begin{frame}{Real-data validation: Co-60}
  \begin{columns}[c]
    \begin{column}{0.48\textwidth}
      \centering
      \shot{user_shots/08_co60_linear.png}{1.0}\\
      \scriptsize Linear
    \end{column}
    \begin{column}{0.48\textwidth}
      \centering
      \shot{user_shots/07_co60_log.png}{1.0}\\
      \scriptsize Logarithmic
    \end{column}
  \end{columns}
\end{frame}

%---------------------------------------------------------------------
\begin{frame}[fragile]{Usage}
  The tool is operated through a menu interface that does not require
  knowledge of ROOT commands.

  \vspace{0.6em}
  \begin{verbatim}
cd UniversalSpectrumToolkit
./ust_menu.sh
  \end{verbatim}

  \vspace{0.4em}
  Menu options:
  \begin{enumerate}
    \item Convert a single file
    \item Batch-convert all files in a directory
    \item Try with sample files
    \item Help
    \item Exit
  \end{enumerate}
\end{frame}

%---------------------------------------------------------------------
\begin{frame}{Menu workflow}
  \begin{center}
    \shot{user_shots/01_menu.png}{0.88}
  \end{center}
\end{frame}

%---------------------------------------------------------------------
\begin{frame}{Decision points in the menu}
  \begin{itemize}
    \item Draw the spectrum on screen: Yes / No
    \item Scale selection: logarithmic Y / linear Y
    \item Forced generic attempt for unknown formats: disabled by default
  \end{itemize}

  \vspace{0.7em}
  \begin{exampleblock}{Design principle}
    The user provides only the file path and visualization preferences.\\
    Format recognition and conversion proceed through validated rules.
  \end{exampleblock}
\end{frame}

%---------------------------------------------------------------------
\begin{frame}{Conclusions}
  \begin{enumerate}
    \item ROOT does not treat \texttt{.spe} content as a native spectrum model.
    \item Reading without a known layout produces physically invalid results.
    \item The solution is resolving the binary layout by hypothesis--score--validation.
    \item The validated layout is encoded as a rule; the output is a ROOT analysis object.
    \item Usage is possible through the menu interface without ROOT expertise.
  \end{enumerate}
\end{frame}

%---------------------------------------------------------------------
\begin{frame}[plain]
  \begin{center}
    \vspace{2.2cm}
    {\Large\bfseries\textcolor{ustnavy}{Thank you for your attention.}}\\[1.4em]
    {\large\textcolor{ustgray}{Respectfully}}\\[1.2em]
    {\normalsize\textcolor{ustteal}{Turkish Accelerator \& Radiation Laboratory / Physics Engineer}}
  \end{center}
\end{frame}

\end{document}
